\scanpage{60}
\mspage{4}
% archive sheet number: 59
et non seulement transitivité à isom. can. près.
Ainsi on peut écrire aussi
\[
\mathrm{HOT}_W(S)\xrightarrow{\sim}\underline{\mathrm{Fib.sc.}}(S)W_S^{-1},\tag*{($\unsure{*|+}$)}
\]
où cette fois $W_S$ est interprété comme
\[
W_S\subset\mathrm{Fl}\bigl(\underline{\mathrm{Fib}}\,\underline{\mathrm{sc}}.(S)\bigr),
\]
savoir
\[
W_S=\{u\in\mathrm{Fl}(-)\mid u_s\in W\quad\substack{\text{pour tt}\\s\in\mathrm{Ob}\,S}\}.
\]
(foncteurs strictement cartésiens, i.e. commutant
aux \unsure{scindages|scindage}, qui induisent des $W$-équivalences
entre les catégories fibres).

La flèche de \unsure{(*)|cas} est bel et bien un iso.
de catégories, pas seulement une équivalence.
Mais on va définir d'autres catégories, plus
commodes à certains égards, qui seront
\uline{équivalentes} (on l'espère), mais sûrement
pas isomorphes, à $\mathrm{HOT}_W(S)$. (On s'occupera
plus tard de la variance de $\mathrm{HOT}_W(S)$ en $S$.)

Les catégories $\underline{\mathrm{Hom}}(S^\circ,\mathrm{Cat})$, $\underline{\mathrm{Fib.sc.}}(S)$,
ainsi que les variantes que nous allons
introduire maintenant, sont en fait des
2-catégories. Mais pour décrire $\mathrm{HOT}_W(S)$,
seules les structures de 1-Catégories sont utilisées. On a l'impression que c'est bien
